\SourceLeaf{249}%
\[
\begin{aligned}
4EEGGk&=E\mathord{.}\frac{\mathrm{d}E}{\mathrm{d}q}\mathord{.}\frac{\mathrm{d}G}{\mathrm{d}q}+E(\frac{\mathrm{d}G}{\mathrm{d}p})^2+G\mathord{.}\frac{\mathrm{d}E}{\mathrm{d}p}\mathord{.}\frac{\mathrm{d}G}{\mathrm{d}p}\\
&\quad+G(\frac{\mathrm{d}E}{\mathrm{d}q})^2-2EG(\frac{\mathrm{d}\mathrm{d}E}{\mathrm{d}q^2}+\frac{\mathrm{d}\mathrm{d}G}{\mathrm{d}p^2})
\end{aligned}
\]
and for the variation of the angle \(\theta\)
\[
\BareRoot{EG}\mathord{.}\mathrm{d}\theta=\tfrac{1}{2}\mathord{.}\frac{\mathrm{d}E}{\mathrm{d}q}\mathord{.}\mathrm{d}p-\tfrac{1}{2}\mathord{.}\frac{\mathrm{d}G}{\mathrm{d}p}\mathord{.}\mathrm{d}q
\]
\par\indent Among the various cases in which this orthogonality condition holds, the most important is the one in which all the lines of either system, for example the first, are shortest lines. Here, for a constant value of \(q\), the angle \(\theta\) is \(=0\). The equation just given for the variation of the angle \(\theta\) therefore shows that we must have \(\frac{\mathrm{d}E}{\mathrm{d}q}=0\), or that the coefficient \(E\) must be independent of \(q\). In other words, \(E\) must be either constant or a function of \(p\) alone. The simplest choice for \(p\) is the length of each line of the first system itself, measured from the point where all the lines of the first system meet whenever there is such a point, or from any line of the second system if there is no common intersection. With these conventions it is clear that \(p\) and \(q\) now denote the same quantities that we represented by \(r\) and \(\varphi\) in Articles 15 and 16, and that we now have \(E=1\). The two preceding formulas thus become the following:
\[
\begin{aligned}
4GGk&=(\frac{\mathrm{d}G}{\mathrm{d}p})^2-2G\frac{\mathrm{d}\mathrm{d}G}{\mathrm{d}p^2}\\
\BareRoot{G}\mathord{.}\mathrm{d}\theta&=-\tfrac{1}{2}\mathord{.}\frac{\mathrm{d}G}{\mathrm{d}p}\mathord{.}\mathrm{d}q
\end{aligned}
\]
or, on setting \(\BareRoot{G}=m\),
\[
k=-\frac{1}{m}\mathord{.}\frac{\mathrm{d}\mathrm{d}m}{\mathrm{d}p^2},\qquad\mathrm{d}\theta=-\frac{\mathrm{d}m}{\mathrm{d}p}\mathord{.}\mathrm{d}q
\]
In general, \(m\) will be a function of \(p,q\), and \(m\mathrm{d}q\) will be the expression for the element of any line of the second system. In the special case where all the lines \(p\) start from the same point, we must clearly have, for \(p=0\), \(m=0\). If in this case we also take for \(q\) the angle that the first element of any line of the first system makes with the element of one of these lines chosen arbitrarily, then for an infinitesimal value of \(p\) the element of the line of the second system, which may be regarded as a circle described with radius \(p\), is \(=p\mathrm{d}q\). We therefore have for an infinitesimal value of \(p,m=p\). Consequently, for \(p=0\), both \(m=0\) and \(\frac{\mathrm{d}m}{\mathrm{d}p}=1\) hold.
\Article{20}
Let us stay with the same assumption a little longer. Let \(p\) denote, in general, the length of a shortest line drawn from a fixed point \(A\) to any point on the surface,
